\documentclass{article}
\usepackage[UTF8]{ctex}
\usepackage{amsmath,amssymb,amsfonts}
\usepackage{enumitem}
\usepackage{geometry}
\geometry{a4paper, margin=1in}
% === 错题本排版优化（不参与题目标识） ===
\setlength{\parindent}{0pt}
\setlength{\parskip}{0.35em plus 0.10em minus 0.08em}
\setlength{\emergencystretch}{3em}
\setlist{itemsep=0.25em, topsep=0.40em, parsep=0pt}
\allowdisplaybreaks[2]
% === 错题本排版优化结束 ===
\begin{document}
\section*{南昌大学2025年数学分析真题}
\section*{120．南昌大学2025年数学分析试卷（回忆版）}

\subsection{题目1}
1． \(\lim _{x \rightarrow 0} \frac{\int_{0}^{x} \sum_{n=0}^{\infty}(-1)^{n} \frac{t^{2 n+1}}{2^{n}(2 n+1)!} \mathrm{d} t-\frac{x^{2}}{2}}{x\left(\sqrt[3]{1+x}+e^{x}\right)}\) \(\_\_\_\_\)．

\subsection{题目2}
2． \(\int \frac{\mathrm{d} x}{1+\sqrt{1-x^{2}}}=\) \(\_\_\_\_\)．

\subsection{题目3}
3．\(\sum_{n=1}^{\infty} \ln \left(\cos \frac{1}{n^{3}}\right) x^{n}\) 收敛半径为 \(\_\_\_\_\)．

\subsection{题目4}
4．\(x+y+z+1=0\) 到椭球 \(18 x^{2}+(\cdot) y^{2}+() z^{2}=()\) 的距离 \(d=\) \(\_\_\_\_\)．

\subsection{题目5}
5．\(u=x^{y} y^{z} z^{x}\) 的全微分 \(\mathrm{d} u=\) \(\_\_\_\_\)．

\subsection{题目6}
6．\(x^{2}+y^{2}=6-z, x z\) 平面，\(y z\) 平面与 \(y=4 z, x=1, y=2\) 围成体积 \(V=\) \(\_\_\_\_\)．

二．（10 分） \(\lim _{n \rightarrow \infty} a_{n}=a,|q|<1\) ，用 \(\varepsilon-N\) 语言证明： \(\lim _{n \rightarrow \infty}\left(a_{n}+a_{n-1} q+\cdots+a_{1} q^{n-1}\right)=\frac{a}{1-q}\) ．

三．（10 分）\(|f(x)-f(y)| \leq M|x-y|\) 为 Lipschitz 条件，\(f(x)\) 满足 \(L\) 条件，\(a>0\) ，证明：\(\frac{f(x)}{x}\) 在\([a,+\infty)\) 上一致连续．

四．（10 分）设 \(f(x)\) 二阶可导，且 \(|f(x)| \leq 1,\left|f^{\prime \prime}(x)\right| \leq 2\) ，证明：\(\left|f^{\prime}(x)\right| \leq 3\) 。（区间可能是 \([0,2]\) ）

五．（10 分）设 \(f(x)\) 在 \([a, b]\) 上无界，证明：对任意的 \(\xi \in[a, b]\) ，存在 \(\delta>0\) ，使得 \(f\) 在 \((\xi-\delta, \xi+\delta) \cap [a, b]\) 上无界．

六．（10分）证明无穷积分的阿贝判别法，即若无穷积分 \(\int_{a}^{+\infty} f(x) \mathrm{d} x\) 收敛，\(g(x)\) 在 \([a,+\infty)\) 上单调有界，则 \(\int_{a}^{+\infty} f(x) g(x) \mathrm{d} x\) 收敛。

七．（10 分）证明：\(f(x)=\left\{\begin{array}{ll}x^{2}, & {[0,1]} \\ 0, & {[0,1]}\end{array}\right.\) 所有有理数；不可积．

八．（15 分）\(f(x)= \begin{cases}\pi-x, & 0<x \leq \pi \\ 0, & x=0 \\ -\pi-x, & -\pi \leq x<0\end{cases}\)\\
（1）写出 \(f(x)\) 的傅里叶级数．\\
（2）证明 \(f(x)\) 在 \((-\pi, \pi)\) 上的一致收敛性（应该是傅里叶级数的一致收敛性）．

九．（15 分）函数项级数 \(\sum_{n=1}^{\infty} n e^{-n x}\) ．\\
（1）证明级数在 \((0,+\infty)\) 上收敛但不一致收敛．\\
（2）求级数的和函数．

\subsection{题目10}
原题附加题：（15 分）求 \(I(\alpha)=\int_{0}^{+\infty} \frac{\arctan x}{x^{\alpha}\left(2+x^{3}\right)} \mathrm{d} x\) 的连续区间．

十一．（15 分）计算

\[
\iint_{\sum}\left(x^{2} \cos (n, x)+y^{2} \cos (n, y)+z^{2} \cos (n, z)\right) \mathrm{d} S .
\]

其中 \(S: x^{2}+y^{2}=z^{2}(0 \leq z \leq h), n\) 为外法向量。
\end{document}
